\section{Rational simplicity and positive rank of the remaining genus-two factors}
\label{qs-section}

This section gives reviewed ordinary proofs, supported by complete finite
tables. It gives no upper rank or complete rational-point classification.
For the fixed branch \(g=3,u=\zeta\), put
\[
 f_t(s)=s^3+3ts^2+(3t-3)s-1,
 \qquad
 H_1:y^2=f_0f_1,\quad H_2:y^2=f_1f_4,\quad H_3:y^2=f_0f_4.
\]
All curves mean their smooth projective models, and
\(J_j=\operatorname{Jac}(H_j)\). These are the three actual quotients
of Section~\ref{hg-section}.

\begin{theorem}[Order-three symmetry and even ranks]\label{qs-symmetry}
Each \(H_j\) has a rational automorphism of exact order three
\[
 \sigma(s,y)=\left(-\frac1{s+1},-\frac{y}{(s+1)^3}\right).
\]
Its quotient has genus zero. There is an embedding
\(\mathbb Q(\sqrt{-3})\hookrightarrow
\operatorname{End}_{\mathbb Q}(J_j)\otimes\mathbb Q\), and
\(\operatorname{rank}J_j(\mathbb Q)\) is even.
\end{theorem}
\begin{proof}
For every \(t\), homogeneous substitution gives
\[
 (s+1)^3 f_t\!\left(-\frac1{s+1}\right)=-f_t(s).
\]
The matrix \(\left(\begin{smallmatrix}0&-1\\1&1\end{smallmatrix}\right)\)
has cube \(-I\), and the weight-three ordinate transformation has the
corresponding sign, so \(\sigma^3=1\). Its action on \(s\) is nontrivial.
The map extends to an automorphism of the smooth projective curve,
including \(s=-1\) and infinity.

Set \(A=-3s(s+1), B=f_1(s), r=A/B\). The two homogeneous cubic
coordinates transform by the same factor. The invariant ordinate
\(y/B\) satisfies, for the three curves respectively,
\[
 v^2=1+r,\qquad w^2=1-3r,\qquad z^2=(1+r)(1-3r).
\]
Since \(s\mapsto r\) has degree three, the extensions of these conic
function fields have degree at most three; the automorphism of order
three forces equality. These are therefore the fixed fields. Each conic
is smooth of genus zero and has the rational point \((r,\text{ordinate})=(0,1)\).
The norm followed by pullback factors through its zero Jacobian and is
\(1+\sigma_*+\sigma_*^2\). Thus \(\sigma_*\) satisfies \(X^2+X+1=0\).
The associated unital map from the quadratic field is injective.
It makes \(J_j(\mathbb Q)\otimes\mathbb Q\) a vector space over that
field, proving that its rational dimension is even, including rank zero.
\end{proof}

\begin{theorem}[Exact certificates at five]\label{qs-five}
The curves \(H_2,H_3\) have good reduction at five, with counts
\[
\begin{array}{c|rr}
 &\#H(\mathbb F_5)&\#H(\mathbb F_{25})\\\hline
 H_2&6&36\\ H_3&6&12
\end{array}
\]
Their Jacobian Frobenius characteristic polynomials are
\begin{equation}\label{qs-five-polynomials}
 P_2(X)=X^4+5X^2+25,\qquad P_3(X)=X^4-7X^2+25.
\end{equation}
\end{theorem}
\begin{proof}
The discriminant of \(f_t\) is \(81(t^2-t+1)^2\), nonzero modulo
five for \(t=0,1,4\). Distinct chosen factors differ by
\(3(t-u)s(s+1)\), with nonzero coefficient, and have values \(-1,1\)
at \(s=0,-1\). Hence their products have six distinct geometric roots
modulo five. Their leading coefficient is one. In the infinity chart
the two points have ordinate \(\pm1\) and are smooth. The two standard
hyperelliptic charts give a smooth proper model over \(\mathbb Z_5\),
and its Jacobian has good reduction.

The sextic values for \(s=0,1,2,3,4\) are respectively
\[
 [1,3,2,3,1],\qquad[1,2,3,2,1].
\]
In both cases the character sum is \(1-1-1-1+1=-1\); adding the
five affine baseline fibres and two points at infinity gives six.
Here is the complete extension-field certificate. In
\(\mathbb F_{25}=\mathbb F_5[\omega]\), \(\omega^2=2\), rows below
index \(b=0,\ldots,4\) and columns \(a=0,\ldots,4\), for inputs
\(a+b\omega\). The entries are quadratic characters of the sextic values:
\[
\begin{array}{c@{\qquad}c}
 H_2&H_3\\[2pt]
 \begin{pmatrix}
 1&1&1&1&1\\1&1&-1&1&1\\1&-1&-1&-1&1\\
 1&-1&-1&-1&1\\1&1&-1&1&1
 \end{pmatrix}
 &
 \begin{pmatrix}
 1&1&1&1&1\\-1&-1&-1&-1&-1\\-1&-1&-1&-1&-1\\
 -1&-1&-1&-1&-1\\-1&-1&-1&-1&-1
 \end{pmatrix}
\end{array}
\]
Their sums are \(9,-15\), giving \(25+9+2=36\) and
\(25-15+2=12\). These entries can be checked using the Legendre
symbol of \(a^2-2b^2\): the norm-character identity is
\(\chi_{25}(x)=x^{12}=N_{25/5}(x)^2\), including the zero convention.
There are no zero entries in these two tables.

For \(S_1=6-\#H(\mathbb F_5)\), \(S_2=26-\#H(\mathbb F_{25})\),
the curve/Jacobian trace formula and Newton's identity give
\[
 P(X)=X^4-S_1X^3+\frac{S_1^2-S_2}{2}X^2-5S_1X+25.
\]
Here \(S_1=0\) and \(S_2=-10,14\), proving
\eqref{qs-five-polynomials}. The standard trace and good-Jacobian
reduction inputs are \cite[Theorem~11.1, Corollaries~11.4 and~12.3]{MilneJV2021}.
\end{proof}

\begin{theorem}[No rational elliptic factors]\label{qs-simple}
Both \(J_2,J_3\) are simple over \(\mathbb Q\) and are not
\(\mathbb Q\)-isogenous to each other. Neither \(H_2\) nor \(H_3\)
has a nonconstant morphism over \(\mathbb Q\) to a genus-one curve.
\end{theorem}
\begin{proof}
If one of these abelian surfaces were not simple over \(\mathbb Q\),
Poincare reducibility over the ground field would give an isogeny to
\(E_a\times E_b\). Its good reduction at five, the rational Tate-module
decomposition and the Neron--Ogg--Shafarevich criterion give good reduction
at five for both elliptic factors; see
\cite[I, Proposition~10.1; IV, Theorem~3.5 and Corollary~3.6]{MilneAVQS}.
Writing their integer traces as \(a,b\) would give
\[
 P_j(X)=(X^2-aX+5)(X^2-bX+5).
\]
The cubic coefficient implies \(b=-a\), and the quadratic coefficient
is then \(10-a^2\). For \(j=2\) one obtains \(a^2=5\); for \(j=3\),
\(a^2=17\). Neither is possible for an integer.

A nonconstant genus-one quotient induces a surjection from \(J_j\)
to a dimension-one Jacobian, contrary to simplicity. One may also use
the rational infinity point on \(H_j\) to identify the target with its
elliptic Jacobian. An isogeny between \(J_2,J_3\) would identify their
rational Tate modules and hence their polynomials at five, which differ.
\end{proof}

This is simplicity over \(\mathbb Q\), not absolute simplicity.
Combining Sections~\ref{hg-section} and~\ref{gd-section} gives
\begin{gather}\label{qs-decomposition}
 \operatorname{Jac}(D_{3,\zeta})\sim_{\mathbb Q}E^2\times J_2\times J_3,
 \\
 \operatorname{rank}\operatorname{Jac}(D_{3,\zeta})(\mathbb Q)
 =2+\operatorname{rank}J_2(\mathbb Q)+\operatorname{rank}J_3(\mathbb Q).
\end{gather}
There are no rational elliptic factors for the two remaining surfaces;
working over an extension field or directly with those surfaces remains
possible. This does not yet determine their ranks or the simultaneous
positive source locus.

\begin{theorem}[Trivial torsion and rank at least six]\label{qs-rank-lower}
Both \(J_2(\mathbb Q)\) and \(J_3(\mathbb Q)\) have trivial torsion
and rank at least two. The genus-six Jacobian in
\eqref{qs-decomposition} therefore has rank at least six.
\end{theorem}
\begin{proof}
The discriminants and factor differences used above also prove good
reduction at seven. The sextic values at \(s=0,\ldots,6\) are
\[
 H_2:[1,0,1,0,3,0,1],\qquad H_3:[1,0,3,0,1,0,1].
\]
Both character sums are two and both curve counts are eleven.
For a complete check over \(\mathbb F_{49}=\mathbb F_7[\omega]\),
\(\omega^2=3\), the character tables, with rows \(b\) and columns
\(a\) in \(a+b\omega\), are
\[
\begin{array}{c@{\quad}c}
 H_2&H_3\\[2pt]
 \begin{pmatrix}
 1&0&1&0&1&0&1\\
 1&-1&1&1&1&-1&1\\
 1&-1&-1&-1&1&1&1\\
 -1&-1&1&-1&1&1&-1\\
 -1&-1&1&-1&1&1&-1\\
 1&-1&-1&-1&1&1&1\\
 1&-1&1&1&1&-1&1
 \end{pmatrix}
 &
 \begin{pmatrix}
 1&0&1&0&1&0&1\\
 1&-1&1&1&1&-1&1\\
 1&1&1&-1&-1&-1&1\\
 -1&1&1&-1&1&-1&-1\\
 -1&1&1&-1&1&-1&-1\\
 1&1&1&-1&-1&-1&1\\
 1&-1&1&1&1&-1&1
 \end{pmatrix}
\end{array}
\]
Both sums are ten, and \(\#H_j(\mathbb F_{49})=49+10+2=61\).
The norm-character rule now uses \(a^2-3b^2\) modulo seven. Newton's
identities give in both cases
\[
 P_7(X)=X^4+3X^3+10X^2+21X+49,\qquad \#J_j(\mathbb F_7)=84.
\]
At five, \eqref{qs-five-polynomials} gives
\(\#J_2(\mathbb F_5)=31\), \(\#J_3(\mathbb F_5)=19\).
Prime-to-\(p\) rational torsion injects into the special fibre at a
good prime \(p\), by the finite etale multiplication kernels and
smooth proper base change (equivalently the specialization used in
\cite[IV, Theorem~3.5]{MilneAVQS}). For a torsion prime other than five
or seven its full primary order must divide both \(84\) and \(31\),
or both \(84\) and \(19\), impossible. Five-primary torsion is excluded
at seven, and seven-primary torsion is excluded at five. All rational
torsion therefore vanishes.

On either \(H_j\), let \(P=(0,1)\) and let \(Q=\infty_+\) have
leading ordinate \(+1\). The class \([P-Q]\) is nonzero. Otherwise a
function with precisely the single simple pole \(Q\) would define a
degree-one map to \(\mathbb P^1\), contrary to genus two. With torsion
zero this class has infinite order. The rank is positive and, by
Theorem~\ref{qs-symmetry}, even, hence at least two. Formula
\eqref{qs-decomposition} completes the proof.
\end{proof}

These are lower bounds. They prevent directly using the classical strict
rank-less-than-genus condition on \(D_{3,\zeta},H_2,H_3\); they do not
exclude quadratic Chabauty, further descent, or methods over extension
fields. No upper rank, Mordell--Weil basis, point-height bound or positive
seed follows from these tables.

\paragraph{Finite verification and ordinary scope.}
The standard-library program
\texttt{replay\_\allowbreak rational\_\allowbreak simple\_\allowbreak
quotients.py --check} regenerates the
complete four curve/prime certificate, directly counting square fibres
and comparing them with norm-character tables. Its canonical SHA-256 is
\texttt{6ba65a743aa19081\allowbreak 9b6d3b43c7e357d0\allowbreak
f5aacb1e53508078\allowbreak cab882a46f240883}.
Both peers independently recomputed the finite counts. The Jacobian,
Frobenius, simplicity and Mordell--Weil conclusions are the ordinary
arguments above with their cited inputs; finite arithmetic is not a
formal proof of the complete geometric conclusions.
